\chapter{Introdução}
\label{cap:introducao}
A pecuária desempenha um papel fundamental na economia brasileira, especialmente nas pequenas e médias propriedades rurais, que representam uma parcela significativa do setor agropecuário nacional. De acordo com o Censo Agropecuário de 2017, existem aproximadamente 5 milhões de estabelecimentos rurais no país, dos quais 73\% são classificados como agricultura familiar de pequenos produtores \cite{centroAgro2017}. Essas propriedades são responsáveis por uma diversidade de produções, incluindo bovinos, suínos, caprinos, equinos e aves, contribuindo significativamente para o abastecimento alimentar do país.

No entanto, a gestão eficiente dessas propriedades apresenta alguns desafios. Muitos produtores ainda dependem de métodos tradicionais e manuais para o gerenciamento de suas atividades, o que pode resultar em falta de controle preciso sobre aspectos essenciais, como histórico dos animais, produção de leite, controle de estoque e registros de vacinação. 

Essa ausência de dados estruturados dificulta a tomada de decisões estratégicas e pode comprometer a sustentabilidade financeira das propriedades.
Embora existam softwares de gestão pecuária disponíveis no mercado, muitos são desenvolvidos com foco em grandes propriedades, apresentando funcionalidades complexas que podem não atender às necessidades específicas das pequenas e médias fazendas. Além disso, a usabilidade desses sistemas nem sempre é intuitiva para produtores que possuem pouca familiaridade com tecnologias digitais. Essa lacuna tecnológica evidencia a necessidade de soluções adaptadas à realidade desses produtores, que sejam acessíveis, fáceis de usar e que considerem a diversidade de animais presentes nas propriedades familiares.


\section{Motivação}
\label{sec:motivacao}

A escolha deste tema decorre de fatores pessoais e profissionais, especialmente considerando a realidade da pecuária no interior do Ceará. Em muitas famílias, mais de 70\% dos membros dependem dessa atividade como principal fonte de renda \cite{centroAgro2017}. No entanto, observa-se a carência de ferramentas tecnológicas acessíveis, intuitivas e adaptadas que organizem de forma estruturada o gerenciamento das propriedades. A falta de sistemas que consolidem informações essenciais, como histórico dos animais, vacinação, produção e estoque, impacta diretamente a eficiência e a sustentabilidade das propriedades. Com experiência de mais de dois anos no desenvolvimento de software, reconhece-se a oportunidade de aplicar conhecimentos técnicos para propor soluções que atendam às necessidades desses pequenos e médios produtores, transformando a gestão rural e potencializando seus resultados.


\section{Objetivos}
Diante desse cenário, este trabalho propõe o desenvolvimento de um aplicativo móvel destinado a auxiliar na gestão de fazendas de pequeno e médio porte, com ênfase na pecuária diversificada. O aplicativo permitirá o mapeamento detalhado do histórico dos animais, monitoramento da produção de leite, controle de estoque de insumos e registro das vacinas administradas em cada animal. Além disso, fornecerá projeções de custos e lucros, visando oferecer aos produtores uma visão clara e estratégica do desempenho de sua propriedade.
\subsection{Objetivo Geral}
\label{sec:objetivo-geral}
O objetivo geral deste trabalho é desenvolver um aplicativo móvel intuitivo e funcional que auxilie na gestão de fazendas de pequeno e médio porte, contemplando diferentes espécies animais. A proposta foca na organização, controle e monitoramento das informações zootécnicas, sanitárias e produtivas, de forma a proporcionar maior eficiência e apoio à tomada de decisões pelos produtores.

\subsection{Objetivos Específicos}
\label{sec:objetivos-especificos}

Para alcançar o supracitado objetivo geral, é possível desdobrá-lo em objetivos específicos a fim de orientar o desenvolvimento deste sistema. Esses objetivos detalham desde uma pesquisa para identificar as reais necessidades do trabalhador pecuário, até o desenvolvimento de funcionalidades. A seguir, os objetivos específicos são melhor detalhados:

	\begin{alineas}
		\item Realizar uma pesquisa para identificar as principais dificuldades enfrentadas pelos proprietários de fazendas em relação à gestão de suas propriedades.
		\item Projetar uma interface de usuário simples e acessível, adequada para produtores com diferentes níveis de familiaridade com tecnologia.
		\item Implementar funcionalidades que permitam o registro e acompanhamento do histórico dos animais, produção de leite, controle de estoque e vacinação.
        \item Desenvolver um módulo de gestão de estoque para insumos, como rações, medicamentos, materiais de manejo e outros recursos utilizados na atividade pecuária.
		\item Desenvolver módulos que forneçam projeções financeiras, auxiliando na tomada de decisões estratégicas.
	\end{alineas}
\section{Estrutura do projeto}
